\subsubsection{Mixing with uniform}
\label{app:mixing}

We also considered a simple data-driven way to mix an adaptive sampling rule with uniform sampling.  Let
\(\pi_m(x)\) denote the sampling probability produced by method \(m\), where \(m\) is either active uncertainty sampling or OPAL, and let \(\pi_{\mathrm{unif}}(x)\) denote the uniform policy with the same expected labeling budget.  For a grid of values \(\lambda \in [0,1]\), we formed
\[
    \pi_{m,\lambda}(x)
    =
    \lambda \pi_m(x) + (1-\lambda)\pi_{\mathrm{unif}}(x),
\]
with probabilities renormalized to preserve the target expected number of labels.  Under this convention, \(\lambda=1\) gives the original adaptive method and \(\lambda=0\) gives pure uniform sampling.

For each budget, we selected \(\lambda\) using the labeled pilot sample.  Specifically, for each candidate value of \(\lambda\), we repeatedly simulated Bernoulli label acquisition from \(\pi_{m,\lambda}\) on the pilot sample and computed the prediction-assisted log odds-ratio estimate.  The tuning criterion was the empirical Monte Carlo variance of this log-scale estimate; equivalently, under the usual Wald approximation this targets the same variance term that determines the confidence interval width.  We then chose
\[
    \lambda_m^*
    =
    \arg\min_{\lambda}
    \widehat{\mathrm{Var}}_{\mathrm{MC}}\{\hat\theta_{\lambda,m}\},
\]
where the expectation is approximated by Monte Carlo simulation over the label-acquisition step.  This tuning rule is intended as a proof of concept: it uses a labeled sample to estimate which amount of uniform exploration would have produced the most stable interval widths at the given budget.

Figure~\ref{fig:sim-lambdatune} shows the selected \(\lambda_m^*\) values across budgets.  OPAL typically selects \(\lambda^*=1\), meaning that the estimated optimum is the unmixed OPAL policy.  In contrast, active uncertainty sampling often selects smaller and budget-dependent values of \(\lambda\), indicating that adding uniform sampling can improve stability.  This behavior is expected because uncertainty sampling can concentrate labels in a narrow region of the score distribution, which can create large inverse-probability weights for observations outside that region.  Uniform mixing raises the low sampling probabilities and can reduce the resulting Monte Carlo variability.

The main limitation is that the selected \(\lambda\) depends on the pilot sample and on the number of Monte Carlo repetitions used in the tuning step.  With a small or unrepresentative pilot sample, the estimated interval widths can be noisy and the chosen \(\lambda\) can overfit the pilot distribution.  Uniform mixing can also be inefficient when the variance-optimal policy is genuinely highly nonuniform, because it spends part of the budget on observations with low influence.  For this reason, we view this experiment primarily as a diagnostic showing that active sampling may require explicit exploration, whereas OPAL is already more stable in this setting.

\begin{figure}[t]
    \centering
    % \includegraphics[width=0.55\linewidth]{figures/uniform_mixing_lambda.pdf}
    \caption{Optimal mixing weight with uniform sampling in the odds-ratio simulation.  For each budget and each base method, \(\lambda^*\) is selected by minimizing the empirical Monte Carlo variance of the prediction-assisted log odds-ratio estimate on the labeled pilot sample.  The convention is \(\lambda=1\) for the unmixed method and \(\lambda=0\) for pure uniform sampling.}
    \label{fig:sim-lambdatune}
\end{figure}
